\documentclass[12pt,a4paper]{article}
\usepackage[utf8]{inputenc}
\usepackage[spanish]{babel}
\usepackage[top=2.5cm,bottom=2.5cm,left=3cm,right=2.5cm]{geometry}
\usepackage{times}
\usepackage{setspace}
\usepackage{graphicx}
\usepackage{float}
\usepackage{booktabs}
\usepackage{array}
\usepackage{longtable}
\usepackage{xcolor}
\usepackage{enumitem}
\usepackage{amsmath}
\usepackage{fancyhdr}
\usepackage{hyperref}
\usepackage{natbib}
\usepackage{listings}
\usepackage{pgfplots}
\usepackage{tikz}
\usetikzlibrary{babel}
\pgfplotsset{compat=1.18}
\onehalfspacing
\setlength{\tabcolsep}{4pt}
\lstset{
  breaklines=true,
  columns=fullflexible,
  keepspaces=true,
  basicstyle=\ttfamily\small
}
\hypersetup{
  colorlinks=true,
  linkcolor=blue,
  citecolor=blue,
  urlcolor=blue
}
\pagestyle{fancy}
\setlength{\headheight}{15pt}
\fancyhf{}
\lhead{UNAP}
\rhead{Sistema de recomendacion preventiva}
\cfoot{\thepage}

\begin{document}
\setcounter{secnumdepth}{0}
\hypersetup{pageanchor=false}
\begin{titlepage}
\centering
{\Large Universidad Nacional del Altiplano\par}
{\large Escuela Profesional de Ingenieria de Sistemas\par}
\vspace{1.2cm}
{\rule{\textwidth}{0.8pt}\par}
\vspace{0.8cm}
{\Huge\bfseries Sistema de Recomendacion de Centros de Salud Preventiva Basado en Capacidad Operativa e Indicadores Climaticos para la Region Puno\par}
\vspace{0.8cm}
{\rule{\textwidth}{0.8pt}\par}
\vspace{1cm}
{\Large Informe tecnico de implementacion funcional\par}
\vspace{1.2cm}
\begin{tabular}{rl}
Unidad academica: & Ingenieria de Sistemas \\
Nivel: & X Semestre \\
Dominio: & Salud publica regional \\
Region: & Puno, Peru \\
Anio: & 2025 \\
\end{tabular}
\vfill
{\large Puno--Peru\par}
\end{titlepage}
\hypersetup{pageanchor=true}
\pagenumbering{roman}
\newpage
\section*{Resumen ejecutivo}
El proyecto desarrollo un sistema funcional de recomendacion para orientar consultas preventivas hacia establecimientos de salud de la region Puno. La solucion integro una API REST en FastAPI, persistencia relacional, autenticacion JWT, un motor de filtrado basado en contenido, una interfaz React con mapa interactivo y un notebook reproducible para analisis exploratorio y evaluacion. La decision de recomendacion combino coincidencia de servicio, categoria IPRESS, distancia geografica, capacidad operativa y penalidad por alerta climatica. El sistema evita datos clinicos individuales y se diseno bajo principios de minimizacion, finalidad y trazabilidad.

\noindent\textbf{Palabras clave:} sistemas de recomendacion, salud preventiva, RENIPRESS, SENAMHI, FastAPI, filtrado basado en contenido, Puno.
\newpage
\tableofcontents
\newpage
\pagenumbering{arabic}

\section{1.- Titulo del proyecto}
\textbf{Sistema de Recomendacion de Centros de Salud Preventiva Basado en Capacidad Operativa e Indicadores Climaticos para la Region Puno.}

El proyecto corresponde al desarrollo de un prototipo funcional de sistema de recomendacion aplicado al dominio de salud publica regional. La propuesta integra datos de establecimientos de salud, atributos operativos, servicios preventivos, ubicacion geografica y alertas climaticas para generar recomendaciones Top-N de centros de atencion preventiva en la region Puno. El sistema fue disenado como una aplicacion web completa, compuesta por backend REST, base de datos relacional, motor de recomendacion, frontend interactivo, notebook analitico e informe tecnico.

El titulo refleja tres decisiones de diseno. Primero, el sistema recomienda \textit{centros de salud preventiva}, por lo que no diagnostica enfermedades ni reemplaza la decision clinica. Segundo, el ranking se basa en \textit{capacidad operativa}, entendida como categoria IPRESS, disponibilidad de servicios, estado operativo y capacidad aproximada. Tercero, incorpora \textit{indicadores climaticos}, debido a que las heladas, precipitaciones intensas y granizadas pueden afectar el traslado de los ciudadanos en zonas altoandinas. En consecuencia, la recomendacion no solo busca similitud entre necesidad e item, sino tambien accesibilidad realista.

\section{2.- Planteamiento del problema}
\subsection{Descripcion del problema}
La region Puno presenta dispersion territorial, zonas rurales de dificil acceso, variacion altitudinal elevada y exposicion frecuente a eventos climaticos adversos. Estas condiciones afectan el acceso oportuno a servicios preventivos como vacunacion, control prenatal, control de crecimiento y desarrollo, tamizaje, nutricion, odontologia y planificacion familiar. Aunque existen establecimientos registrados en RENIPRESS/SUSALUD, la informacion disponible para el ciudadano suele estar distribuida, no siempre se consulta de forma integrada y no incorpora de manera directa el riesgo climatico del momento.

El problema central es la falta de un mecanismo digital que oriente al ciudadano hacia establecimientos preventivos adecuados considerando, al mismo tiempo, oferta de servicios, distancia, categoria del establecimiento, estado operativo y alerta climatica. Sin esa integracion, la eleccion del establecimiento depende de conocimiento local, recomendaciones informales o cercania aparente. Esta situacion puede producir desplazamientos innecesarios, saturacion diferencial de algunos establecimientos, perdida de oportunidad preventiva y mayor exposicion a condiciones climaticas peligrosas.

Desde la perspectiva de sistemas de recomendacion, el desafio consiste en ordenar items disponibles (establecimientos) para un usuario que expresa una necesidad concreta (tipo de atencion preventiva) bajo restricciones geograficas y ambientales. A diferencia de dominios comerciales o de entretenimiento, el error de recomendacion puede afectar tiempo, costo de traslado y seguridad. Por ello, el algoritmo debe ser explicable, auditable y prudente.

\subsection{Contexto}
El contexto del prototipo es la salud publica regional. El sistema esta orientado a ciudadanos de Puno que requieren informacion para acceder a atencion preventiva. El sector salud peruano ha documentado brechas de acceso y continuidad de servicios, especialmente en regiones con barreras geograficas \citep{minsa2023}. La transformacion digital en salud se considera una via para reducir desigualdades si se implementa con criterios de inclusion, interoperabilidad y gobernanza \citep{ops2023}.

Las fuentes de referencia son RENIPRESS/SUSALUD para establecimientos de salud y SENAMHI para informacion meteorologica regional. RENIPRESS permite identificar IPRESS autorizadas y atributos basicos del establecimiento \citep{susalud2026}; SENAMHI publica informacion meteorologica util para valorar condiciones de traslado \citep{senamhi2026}. El prototipo usa estas fuentes como base conceptual y operativa, complementadas con datos reproducibles de ejemplo para ejecucion local.

\subsection{Necesidad del sistema de recomendacion}
La necesidad del sistema se justifica porque el ciudadano no solo requiere una lista de establecimientos, sino un ranking contextualizado. Un listado simple no diferencia adecuadamente entre establecimientos cercanos pero sin el servicio solicitado, establecimientos con buena capacidad pero lejos del usuario, o establecimientos en zonas con alerta climatica activa. Un recomendador permite combinar varios criterios en una unica salida ordenada y explicable.

La tecnica seleccionada fue filtrado basado en contenido porque el proyecto no dispone inicialmente de historiales reales de interaccion. Segun \citet{ricci2022}, este enfoque es apropiado cuando se cuenta con atributos descriptivos de los items y se requiere recomendar sin depender de calificaciones previas. En salud publica, esta decision tambien reduce riesgos de privacidad, ya que no exige historiales clinicos individuales.

\section{3.- Objetivos}
\subsection{Objetivo general}
Desarrollar un sistema de recomendacion funcional que permita orientar a ciudadanos de la region Puno hacia establecimientos de salud preventiva, utilizando filtrado basado en contenido, capacidad operativa, distancia geografica e indicadores climaticos.

\subsection{Objetivos especificos}
\begin{itemize}[leftmargin=*]
  \item Implementar una tecnica de recomendacion basada en contenido para generar Top-N establecimientos preventivos.
  \item Analizar y transformar datos de establecimientos, servicios, coordenadas y alertas climaticas mediante un notebook reproducible.
  \item Desarrollar una API REST segura con autenticacion JWT, persistencia relacional y endpoints de consulta, historial y administracion.
  \item Construir una interfaz web funcional que consuma la API real, permita autenticacion, registre consultas y presente resultados en tarjetas, graficos y mapa interactivo.
  \item Evaluar el sistema mediante metricas de recomendacion, pruebas automatizadas y analisis de riesgos tecnicos.
\end{itemize}

\begin{table}[H]
\centering
\caption{Objetivos especificos y evidencia de implementacion.}
\label{tab:objetivos}
\begin{tabular}{p{5.6cm}p{7.2cm}}
\toprule
Objetivo especifico & Evidencia tecnica \\
\midrule
Integrar datos abiertos de establecimientos y clima & Servicios \texttt{renipress\_service.py} y \texttt{clima\_service.py}; seed reproducible. \\
Implementar recomendacion explicable & Motor \texttt{recommender.py} con similitud del coseno, distancia y penalidad. \\
Exponer API segura & FastAPI, JWT, hashing bcrypt y rutas protegidas. \\
Presentar resultados al ciudadano & React, Leaflet, tarjetas de score, grafico y banner climatico. \\
Documentar y evaluar el modelo & Notebook con EDA, metricas y serializacion. \\
\bottomrule
\end{tabular}
\end{table}

\section{4.- Justificacion}
El sistema es importante porque integra informacion publica dispersa en una herramienta operativa de apoyo a la decision ciudadana. En un entorno altoandino, la distancia y el clima no son variables secundarias: pueden definir si una persona puede trasladarse con seguridad o si conviene priorizar un establecimiento alternativo. Por esa razon, el prototipo aporta valor al combinar datos sanitarios y meteorologicos en un unico ranking.

Los beneficios esperados se agrupan en tres niveles. Para el ciudadano, el sistema reduce incertidumbre al mostrar recomendaciones ordenadas, distancia aproximada, servicios disponibles, horario, score y alerta climatica. Para la gestion publica, el historial de consultas puede servir como insumo agregado para detectar demanda preventiva por territorio, siempre bajo criterios de anonimato y minimizacion. Para el ambito academico, el proyecto demuestra una arquitectura completa de recomendacion desplegable, no limitada a una visualizacion estatica.

La aplicabilidad en el mundo real es alta porque el diseno utiliza tecnologias abiertas y ampliamente usadas: FastAPI, SQLAlchemy, PostgreSQL, React, Leaflet, Pandas y Scikit-learn. Ademas, evita depender de datos clinicos individuales. Esta decision se alinea con la Ley N. 29733 y su reglamento \citep{ley29733,ds0032013}, al reducir la exposicion de datos sensibles y trabajar con datos operativos de establecimientos. La literatura reciente senala que los recomendadores en salud deben priorizar seguridad, explicabilidad y evaluacion contextual, no solo exactitud matematica \citep{sun2023,akbari2023}.

\section{5.- Fases del prototipo}
\subsection{5.1.- Fase 1: Inicial}
\subsubsection{Definicion del dominio}
El dominio seleccionado fue salud publica preventiva en la region Puno. Los servicios considerados incluyen vacunacion, control prenatal, control de crecimiento y desarrollo, tamizaje, odontologia, psicologia, nutricion y planificacion familiar. La eleccion del dominio responde a una necesidad publica concreta: orientar consultas preventivas antes de que el ciudadano realice desplazamientos costosos o riesgosos.

\subsubsection{Identificacion de usuarios, items e interacciones}
Los usuarios del sistema son ciudadanos que requieren atencion preventiva. No se modelan como pacientes clinicos ni se registran diagnosticos. Los items son establecimientos de salud de categorias I-1 a II-2, con atributos como categoria, distrito, provincia, coordenadas, horario, servicios disponibles, estado operativo y capacidad aproximada. El tipo de interaccion principal es una consulta preventiva: el usuario indica ubicacion, edad, tipo de atencion, distancia maxima y numero de recomendaciones.

\begin{table}[H]
\centering
\caption{Elementos del dominio del prototipo.}
\label{tab:dominio}
\begin{tabular}{p{3.5cm}p{9.2cm}}
\toprule
Elemento & Definicion en el sistema \\
\midrule
Usuarios & Ciudadanos de Puno que buscan orientacion para atencion preventiva. \\
Items & Establecimientos de salud registrados o compatibles con RENIPRESS/SUSALUD. \\
Interaccion & Consulta autenticada que genera recomendaciones y queda persistida. \\
Contexto & Ubicacion del usuario, distancia maxima, tipo de atencion y alerta climatica. \\
Salida & Ranking Top-N con score, distancia, servicios, categoria y nivel de alerta. \\
\bottomrule
\end{tabular}
\end{table}

\subsection{5.2.- Fase 2: Analisis exploratorio de datos (EDA)}
\subsubsection{Fuente de datos y descripcion del dataset}
El prototipo utilizo datos compatibles con RENIPRESS/SUSALUD y SENAMHI. Para ejecucion reproducible se incluyeron archivos semilla con establecimientos de Puno, servicios preventivos y alertas climaticas. El dataset contiene variables de identificacion, ubicacion, categoria, servicios, capacidad, horario, estado operativo y nivel de alerta. Esta decision permite ejecutar el proyecto sin depender de disponibilidad externa inmediata, manteniendo una estructura equivalente a las fuentes oficiales.

\begin{longtable}{p{3.4cm}p{4.2cm}p{6cm}}
\caption{Diccionario resumido de datos principales.}
\label{tab:diccionario}\\
\toprule
Entidad & Campo & Uso en el sistema \\
\midrule
\endfirsthead
\toprule
Entidad & Campo & Uso en el sistema \\
\midrule
\endhead
Establecimiento & codigo\_renipress & Identificador unico para trazabilidad. \\
Establecimiento & categoria & Codificacion one-hot del perfil de capacidad. \\
Establecimiento & latitud, longitud & Calculo de distancia haversine. \\
Establecimiento & capacidad\_camas & Aproximacion de capacidad operativa. \\
Servicio & tipo\_servicio & Coincidencia con tipo de atencion requerido. \\
Alerta climatica & nivel\_alerta & Penalidad climatica del score final. \\
Consulta & edad, tipo\_atencion & Vector de usuario y persistencia historica. \\
Recomendacion & score\_final, rank & Ordenamiento Top-N y auditoria. \\
\bottomrule
\end{longtable}

\subsubsection{Limpieza de datos}
La limpieza considero eliminacion de duplicados por codigo RENIPRESS, normalizacion de categorias, normalizacion de nombres de servicios, validacion de coordenadas dentro del rango geografico de Puno e imputacion controlada de campos no criticos como horario. Las coordenadas se trataron como variables criticas porque afectan directamente el calculo de distancia. Los registros con coordenadas fuera de rango deben rechazarse o enviarse a revision antes de ingresar a produccion.

\subsubsection{Transformaciones realizadas}
Las principales transformaciones fueron codificacion one-hot de categoria, binarizacion multietiqueta de servicios, normalizacion de capacidad y calculo de distancia usuario-establecimiento mediante haversine. Para el clima, el nivel de alerta se transformo en penalidad mediante una funcion lineal discreta: nivel 0 equivale a penalidad 0.0, nivel 1 a 0.3, nivel 2 a 0.6 y nivel 3 a 0.9. Esta penalidad reduce el score final, pero no elimina automaticamente el establecimiento; de esa forma, el sistema conserva transparencia y permite evaluar alternativas cuando todas las opciones cercanas estan afectadas por clima adverso.

\subsubsection{Analisis exploratorio}
El notebook genero distribuciones por categoria, provincia, servicios preventivos, matriz de establecimientos por servicios y mapa Folium. Tambien se analizo sparsity de la matriz item-servicio. Este analisis fue necesario porque una matriz excesivamente dispersa puede limitar la discriminacion del recomendador; sin embargo, en el prototipo la combinacion de servicios, categoria, distancia y clima provee suficiente variabilidad para generar rankings diferenciados.

\subsection{5.3.- Fase 3: Implementacion del modelo de recomendacion}
\subsubsection{Tecnica utilizada}
La tecnica utilizada fue filtrado basado en contenido. El sistema representa cada establecimiento mediante atributos observables y compara dicho vector con el perfil de consulta del usuario. La similitud se calcula con coseno, una metrica adecuada para comparar direccion de vectores en espacios donde las magnitudes absolutas pueden variar.

\subsubsection{Algoritmos implementados}
El algoritmo implementado ejecuta las siguientes operaciones: filtra establecimientos activos, calcula distancia haversine, descarta items fuera del radio maximo, obtiene alerta climatica por UBIGEO, construye vectores, calcula similitud del coseno, aplica penalidad climatica, ordena de forma descendente y devuelve Top-N. La formula de similitud fue:
\begin{equation}
\mathrm{sim}(u,i)=\frac{\vec{u}\cdot\vec{i}}{\|\vec{u}\|\|\vec{i}\|}
\label{eq:coseno}
\end{equation}

La penalidad climatica fue:
\begin{equation}
p(n)=0.3n,\quad n\in\{0,1,2,3\}
\label{eq:penalidad}
\end{equation}

El score final fue:
\begin{equation}
score(u,i)=\mathrm{sim}(u,i)\times(1-p(n))
\label{eq:score}
\end{equation}

\begin{lstlisting}[caption={Pipeline logico del motor de recomendacion.},label={lst:pipeline}]
Entrada:
  perfil = {lat, lon, tipo_atencion, edad, max_distancia_km, top_n}

Proceso:
  1. establecimientos = activos()
  2. candidatos = filtrar(haversine(usuario, establecimiento) <= max_distancia)
  3. alerta = alerta_actual_por_ubigeo(candidato.ubigeo)
  4. v_usuario = encode(tipo_atencion, edad, ubicacion, max_distancia)
  5. v_item = encode(categoria, servicios, distancia, capacidad, clima)
  6. similitud = cosine_similarity(v_usuario, v_item)
  7. score_final = similitud * (1 - nivel_alerta * 0.3)
  8. retornar top_n ordenado descendentemente

Salida:
  [{rank, establecimiento, distancia_km, score_final, nivel_alerta, servicios}]
\end{lstlisting}

\subsubsection{Generacion de Top-N recomendaciones}
La salida del modelo es una lista ordenada de establecimientos. Cada recomendacion incluye posicion, nombre, categoria, distancia, score final, nivel de alerta, horario, servicios y coordenadas. El ranking se persiste para permitir auditoria posterior y para formar una base de interacciones que podria alimentar un modelo hibrido en futuras versiones.

\begin{table}[H]
\centering
\caption{Comparacion entre el enfoque implementado y alternativas de recomendacion.}
\label{tab:comparacion_modelos}
\small
\begin{tabular}{@{}p{3cm}p{3.8cm}p{3.4cm}p{2.9cm}@{}}
\toprule
Enfoque & Ventajas & Limitaciones & Pertinencia actual \\
\midrule
Content-Based Filtering & Explicable, no requiere historial, aprovecha atributos RENIPRESS/SENAMHI. & Puede sobrevalorar items similares y depende de calidad de atributos. & Alta. Es el enfoque base implementado. \\
Filtrado colaborativo SVD/NMF & Aprende patrones de usuarios reales y preferencias agregadas. & Requiere historial suficiente y controles de privacidad. & Media futura. Aun no hay datos historicos reales. \\
Modelo hibrido & Combina atributos, historial y reglas de negocio. & Mayor complejidad de despliegue, monitoreo y validacion. & Alta futura cuando exista volumen de uso. \\
Reglas manuales & Facil de auditar y gobernar. & Rigido, baja personalizacion y dificil escalamiento. & Util como capa de seguridad, no como unico motor. \\
\bottomrule
\end{tabular}
\end{table}

\subsection{5.4.- Fase 4: Evaluacion del modelo}
La evaluacion del modelo se abordo desde tres dimensiones: exactitud de ranking, comportamiento operacional y verificacion funcional. Para ranking se calcularon Precision@K y Recall@K con relevancia basada en coincidencia de servicio preventivo. Tambien se incluyo un baseline SVD sobre interacciones sinteticas para mostrar una ruta futura de filtrado colaborativo; este resultado se reporto como exploratorio porque no existe historial real suficiente.

\begin{table}[H]
\centering
\caption{Metricas de evaluacion del enfoque.}
\label{tab:metricas}
\begin{tabular}{lcc}
\toprule
Modelo & Precision@5 & Recall@5 \\
\midrule
Content-Based Filtering & 0.80 & 0.57 \\
Baseline SVD sintetico & No aplica & No aplica \\
\bottomrule
\end{tabular}
\end{table}

La interpretacion de resultados indica que el enfoque basado en contenido es adecuado para una primera version porque produce recomendaciones explicables y no depende de calificaciones historicas. La Precision@5 observada muestra que, en los casos de prueba, la mayoria de elementos recomendados coinciden con el servicio solicitado. El Recall@5 puede mejorar al ampliar el catalogo de establecimientos, enriquecer servicios y calibrar pesos de distancia, capacidad y clima con retroalimentacion real.

\subsection{Herramientas}
El desarrollo utilizo Python, Pandas, Scikit-learn, SQLAlchemy, FastAPI, pytest, React, Vite, Tailwind CSS, React Query, Axios, Leaflet, Recharts, LaTeX y Docker. En el notebook se contemplaron Pandas, Matplotlib, Seaborn, Folium, Scikit-learn, Surprise y joblib.

\section{6.- Desarrollo del sistema}
\subsection{Descripcion tecnica del sistema}
El sistema desarrollado corresponde a una plataforma web de recomendacion operacional orientada a ciudadanos que requieren atencion preventiva en la region Puno. Desde el punto de vista de ingenieria, la solucion se estructuro como una aplicacion cliente-servidor desacoplada: el frontend gestiona la experiencia de usuario, la API concentra reglas de negocio y seguridad, el motor de recomendacion calcula el ranking Top-N y la base de datos conserva entidades persistentes para auditoria, historico de consultas y actualizacion de fuentes externas. Esta separacion permite modificar la interfaz, reemplazar el algoritmo o migrar la base de datos sin reescribir todo el producto.

El dominio del sistema se delimito a establecimientos de salud preventivos y a consultas de orientacion. Por tanto, la recomendacion no constituye diagnostico, prescripcion ni triage clinico. La salida del algoritmo es un ordenamiento de establecimientos segun utilidad operacional esperada: cercania al usuario, disponibilidad del servicio solicitado, categoria del establecimiento, capacidad aproximada y condiciones climaticas asociadas al UBIGEO. Esta definicion resulta importante porque evita atribuir al modelo una competencia medica que no posee y ubica la decision en el plano de accesibilidad y oportunidad.

El flujo de datos inicia con establecimientos compatibles con RENIPRESS/SUSALUD y alertas climaticas compatibles con SENAMHI. La carga inicial se materializa en archivos semilla reproducibles y servicios de sincronizacion. En una instalacion productiva, dichos servicios pueden reemplazarse por conectores directos hacia fuentes oficiales, siempre que se mantenga el mismo contrato de datos. La arquitectura conserva este punto de extension mediante los modulos \texttt{renipress\_service.py} y \texttt{clima\_service.py}.

La capa de recomendacion utiliza filtrado basado en contenido. Cada establecimiento se representa como un vector compuesto por categoria codificada, servicios disponibles, capacidad normalizada y factor climatico. El usuario se representa por el tipo de atencion, edad normalizada, ubicacion y distancia maxima aceptada. El calculo usa similitud del coseno y un ajuste multiplicativo por penalidad climatica. La ventaja de este enfoque es su explicabilidad: cada reduccion del score puede trazarse a una alerta climatica, a la ausencia de un servicio o a una distancia superior al radio de busqueda. En sistemas de salud publica, esta propiedad es preferible a modelos opacos cuando todavia no existe historial real de interacciones.

\subsection{Arquitectura}
\begin{figure}[H]
\centering
\begin{tikzpicture}[node distance=1.5cm, every node/.style={draw, rounded corners, align=center, minimum width=3.2cm, minimum height=0.9cm}]
\node (front) {Frontend\\React + Leaflet};
\node (api) [below of=front] {API REST\\FastAPI + JWT};
\node (engine) [below left of=api, xshift=-2cm] {Motor CBF\\Coseno + clima};
\node (db) [below right of=api, xshift=2cm] {PostgreSQL\\RENIPRESS/SENAMHI};
\node (admin) [below of=api, yshift=-2.8cm] {Admin\\Sync + retrain};
\draw[->] (front) -- (api);
\draw[->] (api) -- (engine);
\draw[->] (api) -- (db);
\draw[->] (admin) -- (db);
\draw[->] (admin) -- (engine);
\end{tikzpicture}
\caption{Arquitectura de componentes del sistema.}
\label{fig:arquitectura}
\end{figure}

El flujo operativo fue: usuario, frontend, API, motor de recomendacion, base de datos y respuesta ordenada Top-N. La arquitectura se dividio en capas para mejorar mantenibilidad: presentacion, servicios de aplicacion, dominio de recomendacion, persistencia y administracion.

\subsection{Componentes clave}
El sistema se organizo en componentes con responsabilidades especificas. Esta descomposicion reduce acoplamiento y facilita pruebas aisladas. La Tabla \ref{tab:componentes} resume el rol de cada modulo principal, su tecnologia y el criterio de calidad asociado.

\small
\begin{longtable}{@{}p{2.8cm}p{2.8cm}p{5.3cm}p{2.4cm}@{}}
\caption{Componentes clave del sistema.}
\label{tab:componentes}\\
\toprule
Componente & Tecnologia & Responsabilidad & Criterio de calidad \\
\midrule
\endfirsthead
\toprule
Componente & Tecnologia & Responsabilidad & Criterio de calidad \\
\midrule
\endhead
Frontend SPA & React, Vite, Tailwind & Gestionar navegacion, formularios, autenticacion en memoria, visualizacion de ranking, mapa y metricas. & Usabilidad, accesibilidad y consumo real de API. \\
Cliente HTTP & Axios & Enviar JWT en cada solicitud protegida y manejar respuestas 401. & Manejo consistente de sesion. \\
API REST & FastAPI & Exponer endpoints versionados, validar schemas, coordinar servicios y persistencia. & Contrato estable y documentado en OpenAPI. \\
Seguridad & python-jose, passlib & Hash de passwords, emision y validacion de tokens JWT. & Confidencialidad de credenciales y control de acceso. \\
Persistencia & SQLAlchemy async & Modelar establecimientos, servicios, alertas, usuarios, consultas y recomendaciones. & Integridad referencial y transacciones. \\
Migraciones & Alembic & Versionar schema relacional para despliegues repetibles. & Reproducibilidad de base vacia. \\
Recomendador & Scikit-learn, NumPy & Construir vectores, calcular similitud, aplicar penalidad y ordenar Top-N. & Exactitud de formula y explicabilidad. \\
Servicios de datos & Pandas, HTTPX & Cargar o actualizar fuentes RENIPRESS/SENAMHI. & Tolerancia a datos incompletos. \\
Planificador & APScheduler & Ejecutar actualizacion climatica periodica. & Vigencia de alertas. \\
Pruebas & pytest, httpx & Verificar autenticacion, ranking, historial y penalidad climatica. & Regresion controlada. \\
\bottomrule
\end{longtable}
\normalsize

\subsection{Flujo de funcionamiento}
El flujo de funcionamiento se diseno para que cada consulta quede autenticada, validada, calculada y persistida. El usuario inicia sesion o se registra; el frontend mantiene el token en memoria y no lo escribe en almacenamiento persistente del navegador. Al enviar una consulta, el cliente remite coordenadas, edad, tipo de atencion, radio maximo y cantidad de resultados. La API valida el esquema con Pydantic, recupera establecimientos activos, calcula distancias, filtra por radio, consulta alertas climaticas actuales y delega el ranking al motor de recomendacion. Finalmente, la consulta y sus recomendaciones se guardan para historial.

\begin{lstlisting}[caption={Diagrama textual del flujo de recomendacion.},label={lst:flujo}]
Ciudadano
  |
  | 1. Login / Registro
  v
Frontend React ---- token JWT en memoria ----+
  |                                          |
  | 2. POST /api/v1/recomendar              |
  v                                          |
API FastAPI -- valida JWT y payload ---------+
  |
  | 3. Consulta establecimientos activos y alertas
  v
PostgreSQL / SQLite local
  |
  | 4. Candidatos dentro del radio
  v
Motor Content-Based
  |
  | 5. cosine_similarity * (1 - penalidad_clima)
  v
Ranking Top-N
  |
  | 6. Persistencia de consulta y recomendaciones
  v
Respuesta JSON + mapa + tarjetas + grafico
\end{lstlisting}

El flujo administrativo se separo del flujo ciudadano. Las rutas administrativas de sincronizacion RENIPRESS, sincronizacion SENAMHI y reentrenamiento requieren un usuario con privilegio \texttt{is\_admin=True}. Esta separacion evita que un usuario comun actualice datos de referencia o reemplace artefactos del modelo. En una version productiva, estos endpoints deberian complementarse con auditoria de cambios, registro de usuario ejecutor, fecha de ejecucion, resumen de registros afectados y resultado de validaciones de calidad de datos.

La persistencia de resultados permite dos beneficios tecnicos. Primero, habilita trazabilidad: es posible reconstruir que establecimientos se recomendaron, con que score y bajo que penalidad. Segundo, crea la base para aprendizaje posterior. Cuando el sistema acumule historiales reales, la tabla \texttt{recomendaciones} puede incorporar retroalimentacion implicita o explicita, por ejemplo visitas confirmadas, clicks en mapa, llamadas o calificaciones. Con esa informacion se podria entrenar un modelo colaborativo o hibrido sin perder el componente explicable basado en contenido.

\subsection{Contrato de API}
La API se organizo bajo el prefijo \texttt{/api/v1}. La Tabla \ref{tab:endpoints} resume los endpoints productivos implementados.

\begin{table}[H]
\centering
\caption{Endpoints REST implementados.}
\label{tab:endpoints}
\begin{tabular}{p{3cm}p{5.2cm}p{5.2cm}}
\toprule
Metodo & Ruta & Responsabilidad \\
\midrule
POST & /auth/register & Crear usuario y emitir token inicial. \\
POST & /auth/login & Validar credenciales y emitir JWT. \\
GET & /establecimientos & Listar establecimientos con filtros. \\
GET & /establecimientos/\{id\} & Obtener detalle y alerta actual. \\
POST & /recomendar & Generar y persistir Top-N recomendado. \\
GET & /historial & Consultar resultados historicos del usuario. \\
GET & /clima/\{ubigeo\} & Obtener alerta climatica actual. \\
POST & /admin/sync-renipress & Sincronizar establecimientos. \\
POST & /admin/sync-senamhi & Sincronizar alertas climaticas. \\
POST & /admin/retrain & Regenerar artefacto del modelo. \\
\bottomrule
\end{tabular}
\end{table}

\subsection{Implementacion en notebook o script}
El notebook contiene la preparacion de datos, EDA, construccion de matriz de items, funcion de recomendacion, ejemplos de consulta, metricas y serializacion de artefactos. El backend replica la logica central como servicio para que el frontend consuma recomendaciones reales mediante la API.

\subsection{Criterios de calidad}
La implementacion se verifico con pruebas unitarias y de integracion. Se evaluaron registro, duplicidad de email, login correcto, password incorrecto, acceso sin token, recomendacion autenticada, penalidad por alerta roja, historial vacio, formula de score y ordenamiento Top-N.

\section{7.- Resultados}
\subsection{Estadisticas del dataset}
\begin{table}[H]
\centering
\caption{Resumen del dataset reproducible de Puno.}
\label{tab:dataset}
\begin{tabular}{lrr}
\toprule
Entidad & Registros & Variables principales \\
\midrule
Establecimientos & 10 & 12 \\
Servicios preventivos & 8 & 2 \\
Alertas climaticas & 10 & 7 \\
Categorias IPRESS & 4 & 1 \\
\bottomrule
\end{tabular}
\end{table}

\subsection{EDA}
El notebook genero distribuciones por categoria, provincia, servicios, matriz de sparsity y mapas Folium. La matriz servicios por establecimiento evidencio dispersion moderada, lo que justifico el uso de codificacion multietiqueta.

\subsection{Ejemplos de recomendacion}
\begin{table}[H]
\centering
\caption{Ejemplo Top-5 para usuario en Puno capital.}
\label{tab:top5}
\begin{tabular}{clccc}
\toprule
Rank & Establecimiento & Categoria & Alerta & Score \\
\midrule
1 & Centro de Salud Metropolitano Puno & I-4 & 1 & 0.61 \\
2 & Puesto de Salud Acora & I-2 & 1 & 0.48 \\
3 & Centro de Salud Ilave & I-4 & 1 & 0.43 \\
4 & Centro de Salud Juliaca Santa Adriana & II-1 & 0 & 0.39 \\
5 & Puesto de Salud Lampa & I-3 & 2 & 0.25 \\
\bottomrule
\end{tabular}
\end{table}

\subsection{Metricas}
La metrica principal del prototipo fue Precision@5, complementada con Recall@5 y pruebas de comportamiento. El resultado no debe interpretarse como desempeno clinico, sino como calidad de ordenamiento respecto al servicio solicitado y restricciones operativas.

\subsection{Metricas de rendimiento y calidad operacional}
Ademas de metricas de ranking, un sistema desplegable requiere observar rendimiento, estabilidad y calidad de servicio. En esta version se verifico funcionalmente con pruebas automatizadas y compilacion de frontend. Para un despliegue productivo se recomienda instrumentar metricas de latencia, tasa de error, uso de CPU, uso de memoria, tiempo de sincronizacion de fuentes, frescura de alertas climaticas y distribucion de scores. Estas metricas no sustituyen las metricas de recomendacion, pero permiten detectar degradacion operacional antes de que impacte al ciudadano.

\begin{table}[H]
\centering
\caption{Metricas recomendadas para monitoreo tecnico.}
\label{tab:metricas_operacionales}
\small
\begin{tabular}{@{}p{3.5cm}p{5.1cm}p{3.8cm}@{}}
\toprule
Metrica & Interpretacion & Umbral sugerido inicial \\
\midrule
Latencia p95 de /recomendar & Tiempo maximo percibido por el 95\% de consultas. & Menor a 800 ms con dataset regional. \\
Tasa de error 5xx & Fallas no controladas del backend. & Menor a 1\% diario. \\
Tasa de 401/403 & Accesos no autorizados o tokens invalidos. & Monitoreo, no necesariamente error. \\
Frescura climatica & Horas desde ultima actualizacion SENAMHI. & Menor a 6 horas. \\
Cobertura de servicios & Porcentaje de establecimientos con cartera preventiva parseada. & Mayor a 95\% tras sincronizacion oficial. \\
Cobertura geoespacial & Porcentaje de establecimientos con coordenadas validas. & Mayor a 98\%. \\
Precision@K & Proporcion de recomendaciones relevantes en K. & Mejorar con datos historicos. \\
Recall@K & Proporcion de relevantes recuperados en K. & Mejorar con catalogo mas completo. \\
Distribucion de penalidad & Frecuencia de resultados por nivel de alerta. & Debe reflejar situacion climatica real. \\
\bottomrule
\end{tabular}
\end{table}

El rendimiento del algoritmo actual es apropiado para un conjunto regional moderado porque opera sobre una matriz de items pequena y calcula similitud vectorial de forma directa. Si el volumen creciera a decenas de miles de establecimientos o si se agregaran redes nacionales, se recomienda introducir cache de matriz de items, precomputo de features estaticos, indices geoespaciales y filtrado previo por bounding box antes de aplicar haversine. Para PostgreSQL, una evolucion natural seria usar PostGIS o, en una primera etapa, indices B-tree por provincia, categoria y estado operativo, combinados con filtros por rango de latitud y longitud.

\begin{table}[H]
\centering
\caption{Pruebas automatizadas y cobertura funcional.}
\label{tab:pruebas}
\small
\begin{tabular}{@{}p{4.2cm}p{6.7cm}p{2.1cm}@{}}
\toprule
Grupo & Casos verificados & Resultado \\
\midrule
Autenticacion & Registro correcto, email duplicado, login correcto, password incorrecto. & Aprobado \\
Recomendacion API & Consulta sin token, consulta autenticada Top-5, penalidad por alerta roja. & Aprobado \\
Historial & Historial vacio de usuario autenticado. & Aprobado \\
Motor unitario & Similitud coseno, penalidad nivel 0, penalidad nivel 3, formula final, orden descendente. & Aprobado \\
Frontend & Compilacion Vite y consumo real de clientes API. & Aprobado \\
Informe & Compilacion LaTeX con BibTeX. & Aprobado \\
\bottomrule
\end{tabular}
\end{table}

\subsection{Sistema funcionando}
La API expuso documentacion interactiva en \texttt{/docs}, autenticacion JWT, endpoints de recomendacion, historial y sincronizacion administrativa. El frontend consumio la API real y mostro mapa Leaflet, marcadores por alerta, grafico de scores y tarjetas Top-N.

\subsection{Despliegue y reproducibilidad}
El backend puede ejecutarse localmente con SQLite para pruebas rapidas o con PostgreSQL 16 mediante Docker Compose. El frontend se compilo con Vite y se sirvio en desarrollo desde el puerto 5173. La documentacion del repositorio incluyo comandos copiables, variables de entorno, carga de datos, ejecucion de pruebas y flujo de uso.

\section{8.- Conclusiones}
El prototipo logro implementar un sistema de recomendacion funcional de extremo a extremo. Se construyo un backend con FastAPI, autenticacion JWT, modelos relacionales, endpoints REST, servicios de sincronizacion, motor de recomendacion y pruebas automatizadas. Tambien se desarrollo un frontend React que consume la API real, presenta resultados en tarjetas, grafico y mapa Leaflet, y permite flujo de login, registro, busqueda, historial y administracion.

Desde el punto de vista del modelo, se implemento filtrado basado en contenido con similitud del coseno y penalidad climatica. Esta tecnica fue adecuada para el escenario inicial porque no requiere historial de usuarios y permite explicar por que un establecimiento fue priorizado o penalizado. La inclusion de distancia y clima hizo que el ranking sea mas realista para el contexto de Puno que una simple busqueda por cercania.

Las principales limitaciones fueron la dependencia de calidad de datos externos, la ausencia de interacciones reales para entrenar filtrado colaborativo y la necesidad de validar en campo la disponibilidad efectiva de servicios. Como aprendizaje principal, se evidencio que en salud publica los recomendadores deben evaluar seguridad, privacidad, equidad territorial y explicabilidad, no solo metricas de exactitud.

\section{9.- Recomendaciones}
Se recomienda ampliar el dataset con registros oficiales actualizados de RENIPRESS y alertas SENAMHI automatizadas. Tambien se recomienda incorporar PostGIS para mejorar consultas geoespaciales, implementar observabilidad con metricas de latencia y frescura climatica, y agregar auditoria administrativa de sincronizaciones y reentrenamientos.

Como extension tecnica, cuando exista historial real de consultas y visitas se podria implementar un modelo hibrido que combine filtrado basado en contenido con SVD, NMF o aprendizaje de ranking. Esta evolucion debe realizarse con controles de privacidad, consentimiento informado, anonimato para analitica y evaluacion de sesgo territorial. En frontend, se recomienda mejorar la explicabilidad visual del score mostrando peso de servicio, distancia, capacidad y clima.

En terminos de seguridad, antes de produccion deben activarse HTTPS, CORS restrictivo, gestion de secretos fuera del repositorio, rate limiting, backups cifrados, politicas de retencion y monitoreo de dependencias. Ademas, el sistema debe comunicar claramente que la recomendacion es orientativa y no reemplaza atencion medica ni confirmacion directa con el establecimiento.

\subsection{Sustento tecnico de las recomendaciones}
El enfoque de datos abiertos evito el uso de historiales clinicos y permitio construir recomendaciones preventivas explicables. Esta estrategia es adecuada para una primera version en salud publica, aunque no sustituye criterios medicos ni decisiones asistenciales. La literatura reciente advierte que los recomendadores en salud deben evaluar calidad, equidad y seguridad, no solo exactitud \citep{vieira2023,sun2023}. Modelos con grafos medicos o datos multimodales ofrecen mejoras cuando existen datos ricos y gobernanza adecuada \citep{deng2024,frontiers2023,tcbb2024}.

Las limitaciones principales fueron la cobertura de estaciones SENAMHI, la posible desactualizacion manual de servicios y el uso de interacciones sinteticas para metricas colaborativas. Cuando exista historial real de consultas y visitas, podra integrarse un modelo hibrido con SVD/NMF, manteniendo reglas de privacidad.

\subsection{Etica, privacidad y gobernanza}
El sistema no expone diagnosticos, tratamientos ni datos clinicos sensibles. La consulta almacena edad, coordenadas aproximadas, tipo de atencion y parametros de busqueda para trazabilidad tecnica. En una puesta en produccion se recomienda incorporar consentimiento informado, retencion limitada, anonimizado para analitica, auditoria de accesos y revision periodica de sesgos territoriales.

\subsection{Analisis de vulnerabilidades y riesgos}
Un sistema de recomendacion aplicado a salud publica debe analizar vulnerabilidades tecnicas y riesgos sociotecnicos. Aunque el sistema no procesa historia clinica ni diagnosticos, maneja datos de ubicacion, edad y patrones de busqueda preventiva. Estos datos pueden revelar condiciones sensibles si se agregan sin controles. Por ejemplo, una consulta de control prenatal desde una coordenada especifica puede inferir informacion personal. Por ello, la proteccion no debe limitarse al cifrado de contrasenas o al JWT; tambien debe incluir minimizacion de datos, retencion limitada, control de acceso, auditoria y anonimizacion para analitica.

\small
\begin{longtable}{@{}p{2.8cm}p{4cm}p{3.1cm}p{3.2cm}@{}}
\caption{Vulnerabilidades potenciales y mitigaciones recomendadas.}
\label{tab:vulnerabilidades}\\
\toprule
Riesgo & Descripcion tecnica & Impacto & Mitigacion \\
\midrule
\endfirsthead
\toprule
Riesgo & Descripcion tecnica & Impacto & Mitigacion \\
\midrule
\endhead
Exposicion de ubicacion & Coordenadas de usuario persistidas con precision completa. & Reidentificacion territorial, especialmente en zonas rurales. & Reducir precision, aplicar geohash truncado, retencion limitada. \\
JWT comprometido & Token robado durante una sesion activa. & Acceso a historial y consultas protegidas. & HTTPS obligatorio, expiracion corta, rotacion de secreto, refresh seguro si se implementa. \\
Secret key debil & Uso de clave por defecto en produccion. & Falsificacion de tokens. & Variables secretas robustas, gestion con vault o secretos del orquestador. \\
Endpoint admin expuesto & Usuario no autorizado intenta sincronizar o reentrenar. & Alteracion de datos o modelo. & Verificacion \texttt{is\_admin}, auditoria, rate limiting, MFA para admin. \\
Datos climaticos obsoletos & Fallo de sincronizacion SENAMHI. & Recomendaciones sin penalidad vigente. & Monitoreo de frescura, alertas operativas, fallback a ultima alerta valida. \\
Calidad deficiente de RENIPRESS & Coordenadas ausentes, servicios incompletos o categoria erronea. & Ranking incorrecto o establecimientos no recomendados. & Validacion de schema, reglas de rango, bitacora de registros rechazados. \\
Sesgo territorial & Zonas con menos datos o estaciones climaticas lejanas reciben peores recomendaciones. & Inequidad de acceso. & Medir cobertura por provincia, incorporar distancia a estacion y revision humana. \\
Ataques de fuerza bruta & Intentos repetidos de login. & Compromiso de cuentas. & Rate limiting, bloqueo temporal, monitoreo de IP, captcha adaptativo. \\
Inyeccion o payload malicioso & Entradas no validadas hacia API. & Error de servicio o abuso. & Pydantic, ORM parametrizado, limites de longitud y tipos estrictos. \\
Dependencias vulnerables & Paquetes npm o pip con CVE. & Compromiso de cadena de suministro. & Auditoria periodica, pinning de versiones, SCA en CI. \\
\bottomrule
\end{longtable}
\normalsize

Desde la perspectiva OWASP, las areas prioritarias son autenticacion, control de acceso, configuracion segura, logging y manejo de dependencias. La implementacion actual ya incorpora hashing de contrasenas y rutas protegidas, pero una puesta en produccion debe agregar controles de infraestructura: HTTPS, cabeceras de seguridad, CORS restringido a dominios oficiales, gestion de secretos fuera del repositorio, backups cifrados, politicas de recuperacion y monitoreo centralizado. En frontend, la decision de almacenar el token solo en memoria reduce persistencia ante ataques XSS, aunque no elimina el riesgo si una sesion activa ejecuta codigo malicioso. Por ello, tambien se recomienda Content Security Policy, sanitizacion de dependencias y revision de cualquier contenido dinamico.

El riesgo algoritimico mas relevante es la sobreconfianza del usuario en el ranking. Un puntaje alto no significa que exista disponibilidad inmediata, personal medico presente o transitabilidad garantizada. El sistema debe comunicar que la recomendacion es una orientacion preventiva basada en datos disponibles. La interfaz ya incorpora alertas climaticas y distancias, pero en produccion deberia incluir fecha de actualizacion de datos, fuente de cada alerta y un indicador de confiabilidad del registro. La transparencia del score es clave: distancia, penalidad climatica y servicios coincidentes deben ser visibles para que el usuario entienda por que un establecimiento fue priorizado.

\subsection{Plan de endurecimiento}
El plan de endurecimiento recomendado se divide en cuatro etapas. La primera etapa corresponde a configuracion segura: cambiar claves por defecto, activar HTTPS, definir CORS estricto y configurar logs estructurados. La segunda etapa corresponde a proteccion de datos: anonimizar coordenadas historicas, definir retencion maxima y separar datos identificables de datos analiticos. La tercera etapa corresponde a observabilidad: metricas de latencia, errores, frescura climatica, calidad de sincronizacion y distribucion territorial de recomendaciones. La cuarta etapa corresponde a gobernanza del modelo: versionado de artefactos, pruebas antes de reentrenar, revision de sesgos y documentacion de cambios.

\begin{lstlisting}[caption={Diagrama textual de controles de seguridad recomendados.},label={lst:seguridad}]
Cliente web
  |-- CSP, no localStorage para JWT, validacion de formularios
  v
API Gateway / Reverse proxy
  |-- HTTPS, rate limiting, CORS estricto, headers de seguridad
  v
FastAPI
  |-- JWT, RBAC admin, Pydantic, logging estructurado
  v
Base de datos
  |-- backups cifrados, roles minimos, retencion limitada
  v
Auditoria y monitoreo
  |-- metricas, trazas, alertas, revision de sesgos
\end{lstlisting}

\clearpage
\renewcommand{\refname}{10.- Referencias}
\addcontentsline{toc}{section}{10.- Referencias}
\bibliographystyle{apalike}
\nocite{gautam2023,healthanalytics2023,ooi2023}
\bibliography{referencias}

\clearpage
\section{11.- Anexos}
\subsection{Anexo A: Codigo fuente}
El codigo fuente se organiza en los directorios \texttt{backend/}, \texttt{frontend/}, \texttt{notebook/} e \texttt{informe/}. El backend contiene la API, modelos ORM, routers, servicios, motor de recomendacion, migraciones, pruebas y archivos de despliegue. El frontend contiene paginas, componentes, hooks, clientes API y configuracion de Vite/Tailwind. El notebook contiene EDA, entrenamiento, evaluacion y serializacion.

\subsection{Anexo B: Dataset}
El dataset reproducible se encuentra en el directorio \texttt{backend/data/}, dentro del archivo \texttt{renipress\_puno\_seed.csv} y archivos asociados de alertas climaticas. Estos datos permiten ejecutar el prototipo localmente y validar el flujo completo sin depender de una conexion externa inmediata.

\subsection{Anexo C: Evidencias adicionales}
Las evidencias tecnicas incluyen documentacion OpenAPI en \texttt{/docs}, pruebas automatizadas con pytest, compilacion del frontend con Vite, generacion del notebook y compilacion LaTeX del informe. Para sustentacion, se recomienda incluir capturas de login, formulario de busqueda, dashboard de resultados, mapa Leaflet, API Swagger y ejecucion de pruebas.
\end{document}
